\scanpage{17}
\mspage{13}
% source: top, date and headings
6. juillet.

\unsure{Re-venons|Revenons}, sur l'introduction des $I_{<a},I_{>b}$.

Autre présentation\unsure{:|.}

% source: upper part, ordered segment
$1^\circ$) Pour y comprendre quelque chose, \unsure{commencons|commençons} avec le cas d'un ``\deleted{ensemble}\inserted{\uline{segment}} \uline{ordonné}'' ou ``\uline{segment orienté combinatoire}'' (peut-être faudrait-il les appeler ``\uline{pseudo-segment}'', car on sera amené par la suite à restreindre les \unsure{hypothèses|hypothéses}), c'est le \uline{système}
\[
\{I,\omega_{\unsure{0|o}},\omega_{\unsure{1|l}}\}\quad\text{ou}\quad(I,\unsure{\partial|\delta}I),\quad\partial I=\{\omega_0,\omega_1\}
\]
d'un ensemble ordonné $I$\inserted{(ou \uline{support du segment orienté})}, d'un plus petit élément $\omega_0$ ou \uline{origine}, et d'un plus grand élément $\omega_1$ ou \uline{extrémité} avec
\[
\omega_0<\omega_1\qquad\text{i.e.}\quad\operatorname{card}I\geq2.
\]

% source: middle, two-element case
On n'exclut pas le cas de $\operatorname{card}I=2$, i.e. $\unsure{\lvert I\rvert^{\circ}|\mathbb I^{\circ}|\lvert I\rvert^0|\mathbb I^0}=\varnothing$, où
\[
\lvert I\rvert^{\circ}=I\setminus\partial I\subset I.
\]
% source: angled left margin
% slanted
\margin{\unsure{C'est le seul} cas où on trouve une maquette}
Ce cas est celui des \uline{segments indivisibles}, qui donnent lieu à une \unsure{m.g. ori} \unsure{qui est} \unsure{une} \uuline{maquette} simpliciale, \inserted{\unsure{savoir},} celle du simplexe orienté-type de dim 1, noté $\Delta^1$.

% source: lower part, restriction to finite intervals
Dans ce cas, on n'introduit pas des strates de la forme $I_{<a},I_{>b}$. La notion du polygone \unsure{complémentaire} ${\unsure{\Phi|\varphi}}^*$ d'un polygone $\Phi$ se développe, en restant dans le cadre des $a$ ou $P_a$ ($a\in I$) et $I_{a,b}$ ($a<b$, $a,b\in I$).

Les développements concernant la détermination
% archive sheet number: 16
